% SPDX-License-Identifier: Apache-2.0
\section{\code{txhdl::types}}
\label{sec:r-types}

Values, and the two markers. \code{Bit} is two valued, \code{Logic} is
IEEE 1164's nine, and \code{U<N>} and \code{I<N>} are the numeric
vectors, with the width a const parameter. \code{I<N>} keeps its bits in
one 128-bit Rust integer, and \code{U<N>} in as many as its second
parameter says, its limbs. \code{U<N>} is \code{U<N, 1>}, one limb, so
its width is at most \code{MAX\_WIDTH}, 128, and a wider one
stops the build where its width is first read, rather than dropping
the bits above. A wider value states its limbs,
\code{U<256, 2>}, since the count cannot be computed from \code{N} on
stable Rust, and the build checks that the count is the one the width
needs. With one limb every operation is the \code{u128} one it always
was; with more, \code{limbs} propagates overflow and borrow bits across limb
boundaries, shifts across them, multiplies in 64-bit digits and divides bit
by bit. What
comes out of a wide value is a word, through \code{slice}, of at most
128 bits; widening into one (\code{concat}, \code{zext},
\code{resize}) and a wide \code{I<N>} are not written yet. A width
that is an
expression of other widths, \code{U<\{A + B\}>}, needs nightly Rust,
which is the largest constraint the embedding found, and the
prototype declares each width by hand instead. The conversions from
Rust integers are what let a literal be written where a value is
expected, and the narrowing ones are checked in debug builds.

\code{Transaction} marks a struct that moves over a channel, and a
bare word is one too. \code{Tag} is a synchronisation domain and its
policy, as a type; which clock edge moves the data is a
\code{Clock}, in the next module.

\lstinputlisting[caption={\code{lib/src/types.rs}.},label={lst:r-types}]{lib/src/types.rs}

\section{\code{txhdl::netlist}}
\label{sec:r-netlist}

The lowering's target, Section~\ref{sec:r-lowering}: \code{Expr},
\code{Target} and \code{Stmt}, which \code{\#[lower]} builds; the
\code{Lowered} unit that holds them with its fields, ports, wires and
memory contents, and, for a unit of units, its nets and its
instances, each a child's \code{Lowered} joined port by port; the
two emitters, \code{verilog} and \code{vhdl}, with the width
inference an extension needs, the hoisting of a slice of an
expression into a wire, and the channel module a net between two
children instantiates, the runtime's buffer of two; the ports sidecar
a testbench generator reads; and, first in the file, the structural
skeleton from the trace walk, which sees fields only.

\lstinputlisting[caption={\code{lib/src/netlist.rs}.},label={lst:r-netlist}]{lib/src/netlist.rs}
